\documentclass[11pt,a4paper]{amsart}
\usepackage[T1]{fontenc}
\usepackage{lmodern}
\usepackage{amssymb}
\usepackage{microtype}
\usepackage[hyphens]{url}
\usepackage{tikz}
\usepackage[colorlinks=true,linkcolor=black,citecolor=black,urlcolor=black]{hyperref}
\usepackage[margin=2.5cm]{geometry}

\hypersetup{
  pdftitle={A short proof that the planes over Q(sqrt3, sqrt11) and Q(sqrt2, sqrt3) are 4-chromatic},
  pdfauthor={Sergi Gal\'an},
  pdfsubject={The Hadwiger-Nelson problem over number fields},
  pdfkeywords={Hadwiger-Nelson problem, chromatic number of the plane, unit-distance graph, number field, formal verification}}

\newtheorem{theorem}{Theorem}
\newtheorem{lemma}[theorem]{Lemma}
\theoremstyle{remark}
\newtheorem*{remark}{Remark}

\newcommand{\F}{\mathbb{F}}
\newcommand{\Q}{\mathbb{Q}}
\newcommand{\R}{\mathbb{R}}
\newcommand{\Z}{\mathbb{Z}}
\newcommand{\p}{\mathfrak{p}}
\newcommand{\Op}{\mathcal{O}_{\mathfrak{p}}}

\title[The planes over $\Q(\sqrt3,\sqrt{11})$ and $\Q(\sqrt2,\sqrt3)$ are $4$-chromatic]{A short proof that the
planes over $\Q(\sqrt3,\sqrt{11})$ and $\Q(\sqrt2,\sqrt3)$ are $4$-chromatic}
\author{Sergi Gal\'an}
\thanks{These results were obtained with the assistance of an AI system (Claude, by Anthropic), which carried
out the computations, the checks and the formalisation in Lean, and helped to draft this note. No mathematician
has checked the note yet. The author is not a professional mathematician and submits it for independent
verification.}
\date{Version 6, 28 September 2026}
\subjclass[2020]{Primary 05C15, 52C10; Secondary 11R16, 68V20}
\keywords{Hadwiger--Nelson problem, chromatic number of the plane, unit-distance graph, number field, formal verification}

\begin{document}

\begin{abstract}
For a subfield $F\subset\R$ let $\chi(F^2)$ be the chromatic number of the graph on $F^2$ joining points at
Euclidean distance $1$. Fischer proved in 1994 that $\chi(\Q(\sqrt3,\sqrt{11})^2)=4$. His result seems to have
been overlooked: Moorhouse (2010), Madore (2015), Cranston--Rabern (2017), Exoo--Ismailescu (2020) and Voronov
(2021) all treat the value as unknown. We give a short proof of a criterion that gives the upper bound $4$ for
this field and for $\Q(\sqrt2,\sqrt3)$; with a $10$-vertex graph it gives $\chi(\Q(\sqrt2,\sqrt3)^2)=4$, the
second case that Voronov thought likely. After this note was written we found that both upper bounds also follow
from a public, unrefereed repository of July 2026 \cite{MM}, whose Theorem~A is the same reduction; it does not
state the case $\Q(\sqrt2,\sqrt3)$. The criterion is a case of
Madore's reduction argument for quadratic forms, in other coordinates. With $\alpha=x+y/\sqrt3$ and
$\beta=2y/\sqrt3$ the squared distance becomes $\alpha^2-\alpha\beta+\beta^2$, a form that is anisotropic modulo
any prime above $2$ with residue field $\F_2$. Both values are also formally verified in Lean~4 with Mathlib.
\end{abstract}

\maketitle

\section{Introduction}

The Hadwiger--Nelson problem asks for $\chi(\R^2)$. Since de Grey \cite{deGrey} and Exoo--Ismailescu
\cite{ExooIsmailescu} it is known that $\chi(\R^2)\in\{5,6,7\}$; see Soifer \cite{Soifer} for the history. For a
subfield $F\subset\R$ one has $\chi(F^2)\le\chi(\R^2)$. For quadratic fields $\Q(\sqrt N)$, with $N$ squarefree:
\begin{itemize}
\item $\chi(\Q^2)=2$ (Woodall \cite{Woodall}), and $\chi(\Q(\sqrt N)^2)=2$ exactly when $N\equiv1,2\pmod 4$
(Johnson \cite{Johnson1987}; Fischer \cite[Thm.~8]{Fischer1990});
\item $\chi(\Q(\sqrt N)^2)\le3$ for $N\equiv0,1\pmod 3$, and $\le4$ for $N\equiv3\pmod 8$ (Fischer
\cite[Thms.~9,~10]{Fischer1990}; see also Payne \cite{Payne});
\item later, Moorhouse \cite{Moorhouse} and Madore \cite{Madore} proved further bounds by reduction modulo a
prime.
\end{itemize}

The field $\Q(\sqrt3,\sqrt{11})$ is the smallest whose plane contains a Moser spindle \cite{Moser} (Moorhouse
\cite[Prop.~1.4]{Moorhouse}). Fischer
\cite{Fischer1994} proved that the plane over $\Q(\sqrt p,\sqrt q)$ has an additive $4$-colouring, with values in
$\Z/4$, for squarefree coprime $p\equiv3\pmod{16}$ and $q\equiv11\pmod{16}$ with $pq\equiv1\pmod{32}$; for
$\Q(\sqrt3,\sqrt{11})$ the Moser spindle gives the lower bound, so $\chi(\Q(\sqrt3,\sqrt{11})^2)=4$. This seems
to have been overlooked. Moorhouse \cite{Moorhouse} noted that the exact value was undetermined, Madore
\cite{Madore} proved $4\le\chi\le5$, Cranston and Rabern \cite{CranstonRabern} asked for the chromatic number,
Exoo and Ismailescu \cite{ExooIsmailescu} asked whether a $5$-chromatic unit-distance graph embeds in
$\Q(\sqrt3,\sqrt{11})^2$, and in Polymath16 Voronov wrote: ``it seems likely that
$\chi(\Q(i,\sqrt3,\sqrt{11}))=4$, and the same is true in case $(2,3)$. But as far as I know, nobody has proved
this yet'' \cite{Polymath17}. Case $(2,3)$ is the field $\Q(i,\sqrt2,\sqrt3)=\Q(\zeta_{24})$, which Fischer's
hypotheses exclude. For it Voronov suggested a colouring through a homomorphism
$\Z[\zeta_{24},1/3]\to(\Z/4)[\zeta_{24}]$, adding ``Perhaps there is a simpler way'' \cite{Polymath17}.

\begin{theorem}\label{thm:main}
$\chi(\Q(\sqrt2,\sqrt3)^2)=\chi(\Q(\sqrt3,\sqrt{11})^2)=4$.
\end{theorem}

The second equality is Fischer's theorem \cite{Fischer1994}; the argument below gives a short proof of both.
Both upper bounds are instances of the following criterion.

\begin{theorem}\label{thm:criterion}
Let $L\subset\R$ be a number field containing $\sqrt3$, and suppose some prime $\p$ of $L$ above $2$ has residue
field $\mathcal{O}_L/\p\cong\F_2$. Then $\chi(L^2)\le4$.
\end{theorem}

Theorem~\ref{thm:criterion} is the case $Q=\alpha^2-\alpha\beta+\beta^2$, with residue field $\F_2$, of Madore's
\P6.6 \cite{Madore}, which extends his Proposition~3.2 to any quadratic form; we include the short proof for
completeness (compare Moorhouse \cite[Lemma~4.2 and Theorem~7.1]{Moorhouse}). In the coordinates $(x,y)$ the
argument fails at $2$: the form $x^2+y^2\equiv(x+y)^2$ is isotropic modulo $2$, and unit vectors such as
$(1/2,\sqrt3/2)$ are not even integral at $2$. Since $\sqrt3\in L$, we have $L(i)=L(\omega)$ with
$\omega=e^{2\pi i/3}$, and the hypothesis on $\p$ implies that $\p$ is inert in $L(i)$: the minimal polynomial
$X^2+X+1$ of $\omega$ has discriminant $-3$, a unit at $\p$, and no root in $\F_2$.

\section{The identity and the lemma}

Put $Q(\alpha,\beta)=\alpha^2-\alpha\beta+\beta^2$. For $x,y\in L$ set
\[
\alpha=x+\frac{y}{\sqrt3},\qquad \beta=\frac{2y}{\sqrt3},\qquad\text{so that}\qquad x^2+y^2=Q(\alpha,\beta).
\]
The identity is a direct expansion. Equivalently, $x+iy=\alpha+\beta\omega$, and $Q$ is the norm form of
$\Q(\omega)$. The map $(x,y)\mapsto(\alpha,\beta)$ is $L$-linear and bijective.

\begin{lemma}\label{lem:valuation}
Let $v$ be the normalised valuation of $L$ at a prime $\p$ above $2$ with residue field $\F_2$. Then for all
$(\alpha,\beta)\ne(0,0)$ in $L^2$,
\[
v(Q(\alpha,\beta))=2\min(v(\alpha),v(\beta)).
\]
\end{lemma}

\begin{proof}
Let $k=\min(v(\alpha),v(\beta))$ and let $\pi$ be a uniformiser. Write $\alpha=\pi^k\alpha_0$ and
$\beta=\pi^k\beta_0$, with $\alpha_0,\beta_0$ integral at $\p$ and not both in $\p$. Then
$Q(\alpha,\beta)=\pi^{2k}Q(\alpha_0,\beta_0)$. Modulo $\p$, $Q(\alpha_0,\beta_0)$ reduces to $s^2+st+t^2$ with
$(s,t)\in\F_2^2\setminus\{0\}$, since signs do not matter in characteristic $2$. This value is $1$ for every
such $(s,t)$, because $X^2+X+1$ has no root in $\F_2$. Hence $v(Q(\alpha_0,\beta_0))=0$.
\end{proof}

\section{Proof of Theorem~\ref{thm:criterion}}

Let $\Op=\{c\in L: v(c)\ge0\}$ and let $\rho\colon\Op\to\F_2$ be reduction modulo $\p$. Fix a set of
representatives of the cosets $L/\Op$, and write $r(c)$ for the representative of $c+\Op$. Since $L$ is
countable, no choice principle is needed. Colour $(x,y)\in L^2$ by
\[
\kappa(x,y)=\bigl(\rho(\alpha-r(\alpha)),\ \rho(\beta-r(\beta))\bigr)\in\F_2^2,
\]
which uses four colours. Suppose $(x,y)$ and $(x',y')$ are at distance $1$, with coordinates $(\alpha,\beta)$
and $(\alpha',\beta')$, and let $\Delta\alpha=\alpha-\alpha'$ and $\Delta\beta=\beta-\beta'$.
\begin{enumerate}
\item By linearity, $Q(\Delta\alpha,\Delta\beta)=1$.
\item By Lemma~\ref{lem:valuation}, $\min(v(\Delta\alpha),v(\Delta\beta))=0$. So $\Delta\alpha$ and
$\Delta\beta$ lie in $\Op$ and are not both in $\p$.
\item Hence $r(\alpha)=r(\alpha')$ and $r(\beta)=r(\beta')$.
\item So the two colours differ by $(\rho(\Delta\alpha),\rho(\Delta\beta))\ne(0,0)$.\qed
\end{enumerate}

\section{Proof of Theorem~\ref{thm:main}}

\emph{Upper bounds.} Both fields satisfy the hypotheses of Theorem~\ref{thm:criterion}.
\begin{itemize}
\item $\Q(\sqrt2,\sqrt3)$. The prime $2$ ramifies in all three quadratic subfields, so no subfield other than
$\Q$ is unramified at $2$, the inertia group is the whole Galois group, and $2=\p^4$ with residue field $\F_2$
($e=4$, $f=g=1$).
\item $\Q(\sqrt3,\sqrt{11})$. The prime $2$ ramifies in $\Q(\sqrt3)$ and $\Q(\sqrt{11})$ ($3,11\equiv3\pmod 4$)
and splits in $\Q(\sqrt{33})$ ($33\equiv1\pmod 8$). So $e\ge2$ and $g\ge2$, and since $efg=4$,
$2=\p_1^2\p_2^2$ and both primes have residue field $\F_2$.
\end{itemize}
Both decompositions were also confirmed with sympy and with PARI/GP (Section~\ref{sec:verification}).

\emph{Lower bounds.} For $\Q(\sqrt2,\sqrt3)$, a $4$-chromatic graph is already implicit in Voronov, Neopryatnaya
and Dergachev \cite{VND}: their second series of $5$-chromatic graphs is built from the $4$-chromatic $10$-vertex
graph $L_{10,2}$ and the unit vectors $\zeta_{24}$ and $(\sqrt6+i\sqrt3)/3$, all with coordinates in
$\Q(\sqrt2,\sqrt3)$. For completeness, here is a short explicit example. Take the unit vectors
\[
u_1=(1,0),\qquad u_2=(0,1),\qquad u_3=\Bigl(-\frac23-\frac{\sqrt2}6,\ -\frac23+\frac{\sqrt2}6\Bigr),
\]
for which $|u_1+u_2+u_3|^2=1/3$. Let $P_0=O$ and $P_k=\sqrt3\,(u_1+\dots+u_k)$. Join $P_{k-1}$ and $P_k$ by the
unit rhombus whose middle vertices are $(P_{k-1}+P_k)/2\pm u_k^\perp/2$, where $u_k^\perp$ is $u_k$ turned by
$90^\circ$. The ten points span $16$ unit distances: the $15$ edges of the three rhombi and $P_0P_3$
(Figure~\ref{fig:chain}).
\begin{enumerate}
\item In any $3$-colouring, the two middle vertices of each rhombus take two distinct colours, and its two tips
both take the third. So each rhombus forces its tips alike, and $P_0$ and $P_3$ share a colour.
\item But $|P_0P_3|=1$, a contradiction.\qed
\end{enumerate}

\begin{figure}[ht]
\centering
\begin{tikzpicture}[scale=2.2]
  % exact coordinates in the comments; the figure is drawn to scale (scripts check the 16 unit distances)
  \coordinate (P0) at (0.0000,0.0000);
  \coordinate (P1) at (1.7321,0.0000);    % (sqrt3, 0)
  \coordinate (P2) at (1.7321,1.7321);    % (sqrt3, sqrt3)
  \coordinate (P3) at (0.1691,0.9856);    % sqrt3 (1/3 - sqrt2/6, 1/3 + sqrt2/6)
  \coordinate (M1p) at (0.8660,0.5000);
  \coordinate (M1m) at (0.8660,-0.5000);
  \coordinate (M2p) at (1.2321,0.8660);
  \coordinate (M2m) at (2.2321,0.8660);
  \coordinate (M3p) at (1.1661,0.9076);
  \coordinate (M3m) at (0.7351,1.8100);
  \draw[line width=0.5pt]
    (P0)--(M1p)--(P1)--(M1m)--(P0) (M1p)--(M1m)
    (P1)--(M2p)--(P2)--(M2m)--(P1) (M2p)--(M2m)
    (P2)--(M3p)--(P3)--(M3m)--(P2) (M3p)--(M3m);
  \draw[line width=1.1pt] (P0)--(P3);
  \foreach \v in {M1p,M1m,M2p,M2m,M3p,M3m} \filldraw[fill=white,draw=black,line width=0.5pt] (\v) circle (0.021);
  \foreach \v in {P0,P1,P2,P3} \fill (\v) circle (0.03);
  \node[below left] at (P0) {$P_0$};
  \node[below right] at (P1) {$P_1$};
  \node[above right] at (P2) {$P_2$};
  \node[left] at (P3) {$P_3$};
\end{tikzpicture}
\caption{The $10$-vertex graph of Section~4, to scale; every edge has length $1$. The rhombi force
$P_0$ and $P_3$ alike in any $3$-colouring, but they are adjacent (thick edge). The white middle vertices of the
last two rhombi are close but distinct.}
\label{fig:chain}
\end{figure}

For $\Q(\sqrt3,\sqrt{11})$ the Moser spindle \cite{Moser} suffices (see Moorhouse \cite[Prop.~1.4]{Moorhouse}): the unit rhombus
with vertices $O$, $(1,0)$, $(1/2,\sqrt3/2)$, $(3/2,\sqrt3/2)$, together with its image under the rotation about
$O$ with $\cos\theta=5/6$, $\sin\theta=\sqrt{11}/6$. The two far tips are at distance $1$, since
$|(3/2,\sqrt3/2)|^2=3$ and $2\cdot3\cdot(1-5/6)=1$.

\section{Remarks}

\begin{enumerate}
\item \emph{Prior work.} Fischer's additive colourings \cite{Fischer1990,Fischer1994} came first. He colours the
plane through the additive group generated by the unit vectors, the component of the origin, whose translates
are the other components \cite[Thm.~1]{Fischer1990}. His $4$-colouring of $\Q(\sqrt N)^2$ for \mbox{$N\equiv3\pmod 8$}
is already a reduction at $2$, with values in $\Z/4$ \cite[Thm.~10]{Fischer1990}, and his $4$-colouring of
$\Q(\sqrt p,\sqrt q)^2$ is also a homomorphism to $\Z/4$ \cite{Fischer1994}. He also described the component of
the origin in the plane over $\Q(\sqrt{N_1},\dots,\sqrt{N_d})$ \cite{Fischer1990b}. Moorhouse
\cite[Lemma~4.2]{Moorhouse} and Madore \mbox{\cite[Prop.~3.2]{Madore}} pass to the whole plane by cosets in the same
way, of a component or of a local ring. Speyer \cite{Polymath2} $4$-coloured the Moser ring
$R=\mathcal{O}_K[1/3]$, $K=\Q(\sqrt{-3},\sqrt{-11})$, in 2018 by reduction modulo $2$:
$R/2R\cong\F_4\times\F_4$, and either projection to $\F_4$ is a $4$-colouring. On $R$ the colouring of
Section~3 is one of these projections, written in the coordinates $\alpha+\beta\omega$. An analysis and computer
search by Hubai, reported by Gibbs, found that all $4$-colourings of $R$ have period $8$ \cite{Polymath3}, and
D\'ucz \cite{Ducz} has recently given geometric $4$-colourings of the Moser lattice and ring. We took the
contribution of this note to be the choice of coordinates, which lets Madore's \P6.6 work at primes above $2$ and
reduces the question to the decomposition of $2$ in $L$, and the case $\Q(\sqrt2,\sqrt3)$. After the note was
written we found the repository \cite{MM} (July 2026, not refereed). Its Theorem~A is the same reduction of
$z=x+iy$ at a conjugation-stable place above $2$, its Corollary~B$'$ gives $\chi\le4$ for a field that contains
$\Q(\sqrt2,\sqrt3)$ and $\Q(\sqrt3,\sqrt{11})$, and it proves $\chi(\Q(\sqrt3,\sqrt{11})^2)=4$. What remains new
here, as far as we know, is the explicit case $\Q(\sqrt2,\sqrt3)$, with its lower bound, and the formal proofs.
That repository also points to a statement for lattices with a condition modulo $8$ on the generators
\cite[Thm.~2.3]{ACLMSS}, which we have not seen.
\item \emph{Consistency check.} The criterion cannot apply to fields that contain $5$-chromatic graphs, and it
does not. In $\Q(\sqrt3,\sqrt{11},\sqrt{247})$, the field of Exoo--Ismailescu, and in
$\Q(\sqrt3,\sqrt5,\sqrt{11})$, the field of Heule \cite{Heule}, every prime above $2$ has residue field
containing $\F_4$, where $X^2+X+1$ has roots.
\end{enumerate}

\section{Verification and provenance}\label{sec:verification}

\begin{itemize}
\item \emph{Upper bounds.} Besides the proof, the colourings were tested by computer on random unit vectors
$z/\bar z$ of each field and on unit-distance graphs with up to $3\,134$ vertices, with no monochromatic edge.
The prime decompositions were recomputed with sympy and with PARI/GP.
\item \emph{Lower bounds.} Both graphs can be checked by hand. An exhaustive computer search also confirms that
neither has a proper $3$-colouring.
\item \emph{Formal proof.} Both equalities of Theorem~\ref{thm:main} are proved in Lean~4 \cite{Lean4}
(version 4.34.1) with Mathlib \cite{Mathlib} (tag v4.34.1): for $L=\Q(\sqrt2,\sqrt3)$ and for
$L=\Q(\sqrt3,\sqrt{11})$, as subfields of $\R$, the graph on $L\times L$ that joins $(x,y)$ and $(x',y')$ when
$(x-x')^2+(y-y')^2=1$ has chromatic number~$4$ (theorems \path{Q23.chromaticNumber_eq_four} and
\path{Q311.chromaticNumber_eq_four}). The proofs use only Lean's standard axioms: propositional
extensionality, quotient soundness and the axiom of choice, which gives the valuation subring and the coset
representatives below. They follow Sections~2--4 in the setting of Madore's \P6.6: $\Op$ is replaced by a
valuation subring $O$ of $L$ with $2$ in its maximal ideal, given by Chevalley's extension theorem;
Lemma~\ref{lem:valuation} is used only in the form ``if $Q(\alpha,\beta)=1$, then $\alpha,\beta\in O$, not both
in the maximal ideal'', which needs no uniformiser; and instead of the decomposition of~$2$, the residue field
of $O$ is shown to be $\F_2$ directly, with an explicit uniformiser for $\Q(\sqrt2,\sqrt3)$ and a Hensel-type
argument for $\Q(\sqrt3,\sqrt{11})$, whichever of its two places above $2$ $O$ belongs to. No ring of integers
is used.
\item \emph{Literature.} We have read \cite{Fischer1990} in full, the zbMATH entries for \cite{Fischer1990b} and
\cite{Fischer1994}, and Johnson's paper \cite{Johnson1987} through its zbMATH entry and Payne's summary
\cite{Payne}. We have not seen the full texts of \cite{Fischer1990b}, \cite{Fischer1994} and
\cite{Johnson1987}, nor Johnson's survey \cite{Johnson2000} of problems in this area.
\item \emph{Code and data.} Everything above is in release 1.1.0 of the repository \cite{Repo}, in the
folder \path{research/hadwiger-nelson}: the formal proofs are \path{lean/LocalColouring.lean},
\path{lean/Q23.lean} and \path{lean/Q311.lean}; the Python tests \path{tests/test_q23.py} and
\path{tests/test_q311.py} check the computations used in the proofs; and the PARI/GP script
\path{scripts/decompositions.gp} recomputes the decompositions of~$2$.
\end{itemize}

\begin{thebibliography}{99}
\bibitem{ACLMSS} M. Axenovich, J. Choi, M. Lastrina, T. McKay, J. Smith and B. Stanton, \emph{On the chromatic
number of subsets of the Euclidean plane}, Graphs Combin. (2014).
\bibitem{CranstonRabern} D. W. Cranston and L. Rabern, \emph{The fractional chromatic number of the plane},
Combinatorica \textbf{37} (2017), 837--861; arXiv:1501.01647.
\bibitem{deGrey} A. D. N. J. de Grey, \emph{The chromatic number of the plane is at least 5}, Geombinatorics
\textbf{28} (2018), 18--31; arXiv:1804.02385.
\bibitem{Lean4} L. de Moura and S. Ullrich, \emph{The Lean 4 theorem prover and programming language},
Automated Deduction -- CADE 28, Lecture Notes in Comput. Sci., vol. 12699, Springer, 2021, pp.~625--635.
\bibitem{Ducz} \'A. D\'ucz, \emph{A note on geometric colorings of the Moser lattice}, arXiv:2606.12325 (2026).
\bibitem{ExooIsmailescu} G. Exoo and D. Ismailescu, \emph{The chromatic number of the plane is at least 5: a new
proof}, Discrete Comput. Geom. \textbf{64} (2020), 216--226; arXiv:1805.00157.
\bibitem{Fischer1990} K. G. Fischer, \emph{Additive $K$-colorable extensions of the rational plane}, Discrete
Math. \textbf{82} (1990), 181--195.
\bibitem{Fischer1990b} K. G. Fischer, \emph{The connected components of the graph
$\Q(\sqrt{N_1},\dots,\sqrt{N_d})^2$}, Congr. Numer. \textbf{72} (1990), 213--221 (Zbl 0733.05048).
\bibitem{Fischer1994} K. G. Fischer, \emph{A planar geometric graph of chromatic number four}, Congr. Numer.
\textbf{104} (1994), 73--79 (Zbl 0836.05030).
\bibitem{Repo} S. Gal\'an, \emph{The Hadwiger--Nelson problem over number fields}, code, data and formal
proofs, release 1.1.0, 2026. Archived at Zenodo,
\href{https://doi.org/10.5281/zenodo.22976635}{doi:10.5281/zenodo.22976635} (all versions); source at
\url{https://github.com/decalion89/chromatic-number-of-the-plane}.
\bibitem{Heule} M. J. H. Heule, \emph{Computing small unit-distance graphs with chromatic number 5},
Geombinatorics \textbf{28} (2018), 32--50; arXiv:1805.12181.
\bibitem{Johnson1987} P. D. Johnson Jr., \emph{Two-colorings of real quadratic extensions of $\Q^2$ that forbid
many distances}, Congr. Numer. \textbf{60} (1987), 51--58 (Zbl 0648.05006).
\bibitem{Johnson2000} P. D. Johnson Jr., \emph{Problems posed in or arising from ``Colorings of metric spaces'':
status report}, Geombinatorics \textbf{9} (2000), no.~4, 170--179.
\bibitem{Madore} D. A. Madore, \emph{The Hadwiger--Nelson problem over certain fields}, arXiv:1509.07023 (2015).
\bibitem{Mathlib} The mathlib Community, \emph{The Lean mathematical library}, Proceedings of the 9th ACM
SIGPLAN International Conference on Certified Programs and Proofs (CPP 2020), ACM, 2020, pp.~367--381.
\bibitem{MM} MildlyMeticulous (GitHub account), \emph{A $2$-adic obstruction to $5$-chromatic unit-distance
graphs, and the exact chromatic number of the plane over $\Q(\sqrt3,\sqrt{11})$}, code, logs and draft, July 2026,
\url{https://github.com/MildlyMeticulous/hn-2adic-obstruction}.
\bibitem{Moorhouse} G. E. Moorhouse, \emph{On the chromatic numbers of planes}, preliminary draft, revised
3 March 2010, \url{https://www.ericmoorhouse.org/pub/chromatic.pdf}.
\bibitem{Moser} L. Moser and W. Moser, \emph{Solution to Problem 10}, Canad. Math. Bull. \textbf{4} (1961),
187--189.
\bibitem{Payne} M. S. Payne, \emph{Unit distance graphs with ambiguous chromatic number}, Electron. J. Combin.
\textbf{16} (2009), Note~31; arXiv:0707.1177.
\bibitem{Polymath2} Polymath16, second thread, \emph{What does it take to be 5-chromatic?}, on D. G. Mixon's
blog \emph{Short, Fat Matrices} (2018), comments 4013--4014 by D. Speyer, 25 April 2018,
\url{https://dustingmixon.wordpress.com/2018/04/22/polymath16-second-thread-what-does-it-take-to-be-5-chromatic/}.
\bibitem{Polymath3} Polymath16, third thread, \emph{Is 6-chromatic within reach?}, ibid.\ (2018), comments 4206
by D. Speyer (3 May 2018) and 4367 by P. Gibbs (9 May 2018),
\url{https://dustingmixon.wordpress.com/2018/05/01/polymath16-third-thread-is-6-chromatic-within-reach/}.
\bibitem{Polymath17} Polymath16, seventeenth thread, \emph{Declaring victory}, ibid.\ (2021), comments 29291 and
29476 by V. Voronov, 18 and 30 July 2021,
\url{https://dustingmixon.wordpress.com/2021/02/01/polymath16-seventeenth-thread-declaring-victory/}.
\bibitem{Soifer} A. Soifer, \emph{The New Mathematical Coloring Book}, Springer, 2024.
\bibitem{VND} V. A. Voronov, A. M. Neopryatnaya and E. A. Dergachev, \emph{Constructing 5-chromatic unit distance
graphs embedded in the Euclidean plane and two-dimensional spheres}, Discrete Math. \textbf{345} (2022), 113106;
arXiv:2106.11824.
\bibitem{Woodall} D. R. Woodall, \emph{Distances realized by sets covering the plane}, J. Combin. Theory Ser. A
\textbf{14} (1973), 187--200.
\end{thebibliography}

\end{document}
